-[6] Lớp 6

----[K] Khoa học tự nhiên

-------[1] Mở đầu về khoa học tự nhiên

----------[1] Giới thiệu về Khoa học tự nhiên
-------------[1] Khái niệm khoa học tự nhiên và đối tượng nghiên cứu
-------------[2] Các lĩnh vực chủ yếu của KHTN (Vật lí, Hoá học, Sinh học, Thiên văn học, Khoa học Trái Đất)
-------------[3] Phân biệt vật sống và vật không sống qua các dấu hiệu đặc trưng
-------------[4] Vai trò của khoa học tự nhiên trong đời sống và sản xuất

----------[2] An toàn trong phòng thực hành
-------------[1] Quy định an toàn và ý nghĩa các kí hiệu cảnh báo trong phòng thực hành
-------------[2] Công dụng và cách dùng dụng cụ đo (thước, bình chia độ, cân, nhiệt kế), kính lúp, kính hiển vi quang học
-------------[3] Xác định giới hạn đo (GHĐ), độ chia nhỏ nhất (ĐCNN) và chọn dụng cụ đo phù hợp
-------------[4] Tính độ phóng đại của kính hiển vi quang học (vật kính x thị kính)
-------------[5] Xử lí tình huống mất an toàn thường gặp trong phòng thực hành

----------[3] Sử dụng kính lúp

----------[4] Sử dụng kính hiển vi quang học

----------[5] Đo chiều dài
-------------[1] Đơn vị và dụng cụ đo chiều dài
-------------[2] Đơn vị và dụng cụ đo khối lượng
-------------[3] Đơn vị và dụng cụ đo thời gian
-------------[4] Thực hành đo kích thước vật dụng trong gia đình (bàn, phòng, sân trường...)
-------------[5] Ước lượng và đo khối lượng thực phẩm khi nấu ăn (cân gạo, đường, muối...)
-------------[6] Giải quyết tình huống thực tế cần đo đạc (mua vải may quần áo, tính diện tích sơn tường...)

----------[6] Đo khối lượng

----------[7] Đo thời gian

----------[8] Đo nhiệt độ
-------------[1] Khái niệm nhiệt độ và thang đo nhiệt độ (°C, °F, K)
-------------[2] Các loại nhiệt kế và cách sử dụng
-------------[3] Thực hành đo nhiệt độ cơ thể, nước, không khí
-------------[4] Ứng dụng đo nhiệt độ trong đời sống (đo sốt, kiểm tra nhiệt độ nước tắm cho trẻ, nhiệt độ bảo quản thực phẩm...)
-------------[5] Đọc hiểu bản tin dự báo thời tiết và ứng dụng vào sinh hoạt hàng ngày

-------[2] Chất quanh ta

----------[1] Sự đa dạng của chất
-------------[1] Nhận biết vật thể tự nhiên, vật thể nhân tạo, vật vô sinh, vật hữu sinh
-------------[2] Đặc điểm các thể cơ bản của chất (Rắn, Lỏng, Khí)
-------------[3] Phân biệt các thể của chất dựa trên tính chất vật lí
-------------[4] Giải thích hiện tượng thực tế liên quan đến thể của chất (Tại sao nước đá nổi trên nước, khí CO₂ chìm xuống...)
-------------[5] Nhận diện và phân loại chất trong các vật dụng gia đình (bàn gỗ, chai nhựa, bóng bay...)

----------[2] Các thể của chất và sự chuyển thể
-------------[1] Sự nóng chảy và sự đông đặc
-------------[2] Sự bay hơi và sự ngưng tụ
-------------[3] Sự sôi
-------------[4] Vẽ và phân tích sơ đồ vòng tuần hoàn của nước trong tự nhiên
-------------[5] Giải thích hiện tượng thực tế về chuyển thể (sương mù buổi sáng, nồi cơm bốc hơi, đá khô trong biểu diễn...)
-------------[6] Ứng dụng sự chuyển thể trong đời sống (đông lạnh bảo quản thực phẩm, sấy khô quần áo, nồi áp suất...)
-------------[7] Phân tích tình huống thực tế liên quan đến nhiệt độ sôi và nhiệt độ nóng chảy (nấu ăn trên núi cao, hàn kim loại...)

----------[3] Oxygen. Không khí
-------------[1] Tính chất vật lí và tầm quan trọng của Oxygen
-------------[2] Sự cháy và các yếu tố cần thiết cho sự cháy
-------------[3] Biện pháp dập tắt đám cháy
-------------[4] Ứng dụng Oxygen trong y tế và đời sống (bình thở, nuôi trồng thủy sản, xử lý nước thải...)
-------------[5] Giải quyết tình huống phòng cháy chữa cháy tại gia đình và trường học
-------------[6] Giải thích hiện tượng thực tế liên quan đến sự cháy (tại sao thổi nến tắt, lửa cháy trong rừng lan nhanh...)
-------------[7] Thành phần của không khí
-------------[8] Thí nghiệm xác định thành phần phần trăm thể tích Oxygen trong không khí
-------------[9] Vai trò của không khí đối với tự nhiên và sự sống
-------------[A] Nguyên nhân và hậu quả của ô nhiễm không khí
-------------[B] Biện pháp bảo vệ môi trường không khí
-------------[C] Đánh giá chất lượng không khí tại địa phương (đọc chỉ số AQI, nhận diện nguồn ô nhiễm quanh nhà/trường)
-------------[D] Đề xuất giải pháp giảm ô nhiễm không khí trong sinh hoạt gia đình (hạn chế đốt rác, trồng cây xanh, sử dụng phương tiện xanh...)
-------------[E] Phân tích tác động của ô nhiễm không khí đến sức khỏe con người (bệnh hô hấp, dị ứng, biến đổi khí hậu...)

-------[3] Một số vật liệu, nguyên liệu, nhiên liệu, lương thực – thực phẩm thông dụng

----------[1] Một số vật liệu
-------------[1] Tính chất và ứng dụng của vật liệu (Kim loại, Nhựa, Gỗ, Thủy tinh, Gốm sứ, Cao su...)
-------------[2] Dấu hiệu nhận biết vật liệu bị hư hỏng
-------------[3] Quy tắc an toàn và hiệu quả khi sử dụng vật liệu (3R)
-------------[4] Lựa chọn vật liệu phù hợp cho mục đích sử dụng cụ thể (xây nhà, làm đồ dùng, may mặc...)
-------------[5] Đánh giá tác động môi trường của các loại vật liệu (nhựa dùng một lần, túi ni-lông, vật liệu tái chế...)
-------------[6] Thực hành phân loại rác và tái chế vật liệu tại gia đình, trường học

----------[2] Một số nguyên liệu
-------------[1] Khái niệm và phân loại nguyên liệu (Quặng, Đá vôi...)
-------------[2] Tính chất và ứng dụng của nguyên liệu trong sản xuất
-------------[3] Khai thác nhiên liệu và phát triển bền vững
-------------[4] Tìm hiểu quy trình sản xuất từ nguyên liệu đến sản phẩm thực tế (đá vôi → xi măng, quặng sắt → thép...)
-------------[5] Đánh giá tác động của khai thác nguyên liệu đến môi trường địa phương (khai thác mỏ, phá rừng...)

----------[3] Một số nhiên liệu
-------------[1] Khái niệm và phân loại nhiên liệu (Rắn, Lỏng, Khí)
-------------[2] Tính chất và ứng dụng của một số nhiên liệu phổ biến (Than, Xăng, Dầu, Gas...)
-------------[3] An ninh năng lượng và sử dụng nhiên liệu an toàn, tiết kiệm
-------------[4] So sánh ưu nhược điểm của các loại nhiên liệu trong sinh hoạt (gas vs bếp từ vs bếp củi)
-------------[5] Tính toán chi phí sử dụng nhiên liệu trong gia đình (tiền gas, tiền xăng xe hàng tháng...)
-------------[6] Đề xuất giải pháp tiết kiệm năng lượng và sử dụng năng lượng tái tạo (pin mặt trời, biogas...)

----------[4] Một số lương thực, thực phẩm
-------------[1] Vai trò của lương thực, thực phẩm và các nhóm chất dinh dưỡng
-------------[2] Tính chất của lương thực, thực phẩm (trạng thái, màu sắc, mùi vị...)
-------------[3] Dấu hiệu ngộ độc thực phẩm và cách phòng tránh
-------------[4] Bảo quản lương thực, thực phẩm đúng cách
-------------[5] Đọc hiểu nhãn thực phẩm (thành phần, hạn sử dụng, chỉ tiêu dinh dưỡng, phụ gia...)
-------------[6] Lập thực đơn cân bằng dinh dưỡng cho bữa ăn gia đình
-------------[7] Nhận diện thực phẩm bẩn, thực phẩm không an toàn trong thực tế (rau có thuốc trừ sâu, thịt ôi thiu, nước giải khát có chất bảo quản...)
-------------[8] Xử lý tình huống ngộ độc thực phẩm (sơ cứu, gọi cấp cứu, báo cáo cơ quan chức năng...)

-------[4] Hỗn hợp. Tách chất ra khỏi hỗn hợp

----------[1] Hỗn hợp các chất
-------------[1] Khái niệm và cách nhận biết chất tinh khiết, hỗn hợp
-------------[2] Phân biệt hỗn hợp đồng nhất và hỗn hợp không đồng nhất
-------------[3] Nhận diện chất tinh khiết và hỗn hợp trong đời sống hàng ngày (nước khoáng, không khí, nước muối, vàng 24K...)
-------------[4] Giải thích tại sao nước máy không phải chất tinh khiết và quy trình xử lý nước sinh hoạt
-------------[5] Khái niệm và phân biệt dung dịch, huyền phù, nhũ tương
-------------[6] Thành phần của dung dịch (Dung môi, Chất tan)
-------------[7] Khả năng hòa tan của chất rắn trong nước (Ảnh hưởng của nhiệt độ, khuấy trộn...)
-------------[8] Tính toán nồng độ phần trăm, tính lượng chất tan/dung môi (Mức độ nâng cao)
-------------[9] Nhận diện dung dịch, huyền phù, nhũ tương trong đời sống (nước chanh, sữa, nước sông, sốt mayonnaise...)
-------------[A] Ứng dụng hòa tan trong nấu ăn và pha chế (pha nước muối súc miệng, pha sữa, pha thuốc...)
-------------[B] Giải thích hiện tượng thực tế về độ tan (tại sao đường tan nhanh hơn trong nước nóng, muối biển kết tinh khi phơi nắng...)

----------[2] Tách chất khỏi hỗn hợp
-------------[1] Phương pháp Lắng, Gạn, Lọc (Nguyên tắc và ứng dụng)
-------------[2] Phương pháp Cô cạn (Nguyên tắc và ứng dụng)
-------------[3] Phương pháp Chiết (Nguyên tắc và ứng dụng)
-------------[4] Lựa chọn phương pháp tách biệt phù hợp cho hỗn hợp cụ thể
-------------[5] Thực hành tách chất (Tách muối ăn, tách dầu ăn khỏi nước...)
-------------[6] Ứng dụng phương pháp tách chất trong đời sống (lọc nước uống bằng bình lọc, làm muối từ nước biển, ép dầu từ lạc/vừng...)
-------------[7] Phân tích quy trình xử lý nước thải sinh hoạt (lắng → lọc → khử trùng)
-------------[8] Giải quyết tình huống thực tế: đề xuất cách tách hỗn hợp gặp trong sinh hoạt (tách cát khỏi gạo, lọc nước mưa để sử dụng...)

-------[5] Tế bào

----------[1] Tế bào – Đơn vị cơ bản của sự sống
-------------[1] Khái niệm tế bào - đơn vị cơ sở của sự sống
-------------[2] Cấu tạo tế bào (màng tế bào, chất tế bào, nhân tế bào)
-------------[3] Phân biệt tế bào nhân sơ và tế bào nhân thực
-------------[4] Phân biệt tế bào động vật và tế bào thực vật
-------------[5] Liên hệ thực tế: quan sát tế bào qua kính hiển vi (tế bào hành tây, tế bào máu...)
-------------[6] Giải thích hiện tượng thực tế liên quan đến tế bào (vết thương lành, da bong tróc, tóc mọc dài...)

----------[2] Cấu tạo và chức năng các thành phần của tế bào

----------[3] Sự lớn lên và sinh sản của tế bào
-------------[1] Sự lớn lên của tế bào
-------------[2] Sự phân chia (sinh sản) của tế bào
-------------[3] Ý nghĩa của sự lớn lên và sinh sản tế bào đối với cơ thể
-------------[4] Giải thích tại sao trẻ em lớn nhanh, vết thương tự lành, móng tay mọc dài
-------------[5] Liên hệ ứng dụng y học: ghép da, nuôi cấy mô trong điều trị bệnh

----------[4] Thực hành: Quan sát và phân biệt một số loại tế bào

-------[6] Từ tế bào đến cơ thể

----------[1] Cơ thể sinh vật
-------------[1] Đặc điểm cơ thể đơn bào (trùng roi, vi khuẩn...)
-------------[2] Đặc điểm cơ thể đa bào (thực vật, động vật, con người)
-------------[3] So sánh cơ thể đơn bào và đa bào
-------------[4] Liên hệ thực tế: vi khuẩn đơn bào gây bệnh trong đời sống (tiêu chảy, ngộ độc thực phẩm...)
-------------[5] Ứng dụng sinh vật đơn bào trong sản xuất (nấm men làm bánh, vi khuẩn lên men sữa chua...)

----------[2] Tổ chức cơ thể đa bào
-------------[1] Khái niệm mô, cơ quan, hệ cơ quan, cơ thể
-------------[2] Mối quan hệ giữa các cấp độ tổ chức
-------------[3] Nhận diện các loại mô, cơ quan trong cơ thể người (mô cơ, mô xương, hệ tiêu hóa...)
-------------[4] Giải thích tại sao khi một cơ quan bị tổn thương ảnh hưởng đến cả cơ thể (đau dạ dày → mệt mỏi toàn thân...)

----------[3] Thực hành: Quan sát và mô tả cơ thể đơn bào, cơ thể đa bào

-------[7] Đa dạng thế giới sống

----------[1] Hệ thống phân loại sinh vật
-------------[1] Sự cần thiết của việc phân loại thế giới sống
-------------[2] Các bậc phân loại sinh vật (Loài, Chi, Họ, Bộ, Lớp, Ngành, Giới)
-------------[3] Sử dụng khóa lưỡng phân để phân loại sinh vật
-------------[4] Thực hành phân loại động vật, thực vật trong vườn trường và khu dân cư
-------------[5] Ứng dụng phân loại sinh vật trong nông nghiệp (nhận diện cây trồng, giống vật nuôi...)

----------[2] Khoá lưỡng phân

----------[3] Vi khuẩn
-------------[1] Đặc điểm và cấu tạo vi khuẩn
-------------[2] Vai trò và tác hại của vi khuẩn
-------------[3] Giải thích hiện tượng thức ăn bị thiu, sữa chua lên men, đất màu mỡ nhờ vi khuẩn
-------------[4] Biện pháp phòng tránh bệnh do vi khuẩn gây ra (rửa tay, vệ sinh an toàn thực phẩm...)
-------------[5] Ứng dụng vi khuẩn có lợi (làm sữa chua, muối dưa, xử lý nước thải, sản xuất phân vi sinh...)

----------[4] Thực hành: Làm sữa chua và quan sát vi khuẩn

----------[5] Virus
-------------[1] Đặc điểm và cấu tạo virus
-------------[2] Một số bệnh do virus gây ra
-------------[3] Biện pháp phòng tránh bệnh do virus (tiêm vaccine, đeo khẩu trang, rửa tay...)
-------------[4] Liên hệ thực tế: phòng chống dịch bệnh (COVID-19, cúm mùa, sốt xuất huyết, sởi...)
-------------[5] Đánh giá vai trò của tiêm chủng trong bảo vệ sức khỏe cộng đồng

----------[6] Nguyên sinh vật
-------------[1] Đặc điểm và phân loại nguyên sinh vật (tảo, trùng...)
-------------[2] Vai trò và tác hại của nguyên sinh vật
-------------[3] Giải thích hiện tượng nước ao hồ có màu xanh (tảo nở hoa)
-------------[4] Phòng tránh bệnh do nguyên sinh vật gây ra (sốt rét, kiết lị - diệt muỗi, ăn chín uống sôi...)

----------[7] Thực hành: Quan sát nguyên sinh vật

----------[8] Nấm
-------------[1] Đặc điểm và phân loại nấm (nấm đơn bào, nấm đa bào)
-------------[2] Vai trò và tác hại của nấm
-------------[3] Phân biệt nấm ăn được và nấm độc trong tự nhiên (nấm rơm, nấm hương vs nấm độc tán trắng...)
-------------[4] Ứng dụng nấm trong đời sống (nấm men làm bánh mì, nấm mốc sản xuất kháng sinh penicillin...)
-------------[5] Phòng tránh bệnh nấm da, nấm móng trong sinh hoạt hàng ngày

----------[9] Thực hành: Quan sát các loại nấm

----------[A] Thực vật
-------------[1] Đặc điểm chung và phân loại thực vật (Rêu, Dương xỉ, Hạt trần, Hạt kín)
-------------[2] Vai trò của thực vật trong tự nhiên và đời sống
-------------[3] Nhận diện và phân loại thực vật tại địa phương (cây lương thực, cây thuốc, cây cảnh...)
-------------[4] Đề xuất giải pháp bảo vệ đa dạng thực vật (trồng cây xanh, bảo tồn rừng, chống phá rừng...)
-------------[5] Ứng dụng thực vật trong y học cổ truyền và đời sống (cây thuốc nam quanh nhà...)
-------------[6] Thiết kế vườn rau/vườn thuốc nam tại gia đình hoặc trường học

----------[B] Thực hành: Quan sát và phân biệt một số nhóm thực vật

----------[C] Động vật
-------------[1] Đặc điểm chung và phân loại động vật (Động vật không xương sống, Động vật có xương sống)
-------------[2] Vai trò của động vật trong tự nhiên và đời sống
-------------[3] Nhận diện động vật thường gặp và phân loại (côn trùng, cá, chim, thú...)
-------------[4] Đề xuất biện pháp bảo vệ động vật hoang dã và đa dạng sinh học tại địa phương
-------------[5] Ứng dụng kiến thức động vật trong chăn nuôi gia đình (nuôi gà, cá, ong...)
-------------[6] Phân tích chuỗi thức ăn và lưới thức ăn trong hệ sinh thái quanh nhà/trường

----------[D] Thực hành: Quan sát và nhận biết một số nhóm động vật ngoài thiên nhiên

----------[E] Đa dạng sinh học

----------[F] Tìm hiểu sinh vật ngoài thiên nhiên

-------[8] Lực trong đời sống

----------[1] Lực là gì?
-------------[1] Khái niệm lực, biểu diễn lực bằng mũi tên
-------------[2] Tác dụng của lực (làm vật biến dạng, thay đổi chuyển động)
-------------[3] Lực tiếp xúc và lực không tiếp xúc
-------------[4] Nhận diện các lực tác dụng lên vật trong đời sống (đẩy xe, kéo cửa, nam châm hút sắt...)
-------------[5] Giải thích hiện tượng thực tế bằng kiến thức về lực (tại sao bóng đổi hướng khi đá, ô tô phanh gấp...)

----------[2] Biểu diễn lực

----------[3] Biến dạng của lò xo

----------[4] Trọng lượng, lực hấp dẫn
-------------[1] Khái niệm lực hấp dẫn, trọng lực
-------------[2] Trọng lượng và cách đo trọng lượng (lực kế)
-------------[3] Giải thích tại sao vật rơi xuống đất, nước chảy từ cao xuống thấp
-------------[4] Liên hệ thực tế: trạng thái không trọng lượng trong vũ trụ, tại sao phi hành gia lơ lửng
-------------[5] Ứng dụng trọng lực trong đời sống (thủy điện, hệ thống tưới tự chảy...)

----------[5] Lực ma sát
-------------[1] Khái niệm lực ma sát (ma sát trượt, ma sát lăn, ma sát nghỉ)
-------------[2] Các yếu tố ảnh hưởng đến lực ma sát
-------------[3] Lực ma sát có lợi và có hại
-------------[4] Giải thích hiện tượng thực tế (tại sao đi giày có đế gai không bị trượt, xe phanh trên đường ướt khó dừng...)
-------------[5] Ứng dụng tăng/giảm ma sát trong đời sống (tra dầu mỡ máy, làm nhám đế giày, phanh xe đạp...)
-------------[6] Phân tích an toàn giao thông liên quan đến ma sát (khoảng cách phanh, lốp xe mòn...)

----------[6] Lực cản của nước
-------------[1] Khái niệm lực cản của nước và không khí
-------------[2] Ảnh hưởng của hình dạng vật đến lực cản
-------------[3] Giải thích thiết kế khí động học (xe hơi, tàu cao tốc, máy bay mũi nhọn...)
-------------[4] Ứng dụng và hạn chế lực cản (dù nhảy, tàu ngầm, áo phao...)
-------------[5] Giải thích tại sao bơi lội mệt hơn đi bộ, đạp xe ngược gió khó hơn xuôi gió

-------[9] Năng lượng

----------[1] Năng lượng và sự truyền năng lượng

----------[2] Một số dạng năng lượng
-------------[1] Nhận biết các dạng năng lượng (động năng, thế năng, nhiệt năng, quang năng, điện năng, hóa năng...)
-------------[2] Năng lượng hao phí và năng lượng có ích
-------------[3] Nhận diện các dạng năng lượng trong sinh hoạt hàng ngày (bếp gas → nhiệt năng, đèn → quang năng, quạt → động năng...)
-------------[4] Phân tích năng lượng trong chuỗi hoạt động thực tế (pin → điện thoại → ánh sáng + âm thanh + nhiệt...)

----------[3] Sự chuyển hoá năng lượng
-------------[1] Khái niệm sự chuyển hóa năng lượng
-------------[2] Định luật bảo toàn năng lượng (mức độ đơn giản)
-------------[3] Nhận diện sự chuyển hóa năng lượng trong các thiết bị gia đình (nồi cơm điện, máy giặt, tivi...)
-------------[4] Giải thích hiện tượng: tại sao điện thoại nóng khi sạc, tay ấm khi xoa vào nhau
-------------[5] Phân tích chuỗi chuyển hóa năng lượng từ nhà máy điện đến thiết bị trong nhà

----------[4] Năng lượng hao phí

----------[5] Năng lượng tái tạo

----------[6] Tiết kiệm năng lượng
-------------[1] Nguồn năng lượng tái tạo và không tái tạo
-------------[2] Sự cần thiết phải tiết kiệm năng lượng
-------------[3] Biện pháp tiết kiệm năng lượng trong gia đình (tắt đèn, dùng thiết bị tiết kiệm điện...)
-------------[4] Tính toán hóa đơn tiền điện gia đình và đề xuất giải pháp giảm chi phí
-------------[5] Đánh giá hiệu quả sử dụng năng lượng tái tạo tại địa phương (pin mặt trời, điện gió, biogas...)
-------------[6] Thiết kế poster/khẩu hiệu tuyên truyền tiết kiệm năng lượng tại trường học

-------[A] Trái Đất và bầu trời

----------[1] Chuyển động nhìn thấy của Mặt Trời. Thiên thể
-------------[1] Chuyển động biểu kiến của Mặt Trời (mọc - lặn)
-------------[2] Các hình dạng nhìn thấy của Mặt Trăng (các pha Trăng)
-------------[3] Giải thích hiện tượng ngày - đêm, mùa trong năm
-------------[4] Ứng dụng hiểu biết về Mặt Trời trong đời sống (chọn hướng nhà, phơi đồ, lắp pin mặt trời...)
-------------[5] Sử dụng lịch âm để biết ngày trăng tròn, trăng non (ứng dụng trong nông nghiệp, lễ hội...)

----------[2] Mặt Trăng

----------[3] Hệ Mặt Trời
-------------[1] Các thành phần của hệ Mặt Trời (8 hành tinh, tiểu hành tinh, sao chổi...)
-------------[2] Đặc điểm của Trái Đất trong hệ Mặt Trời
-------------[3] Khái niệm cơ bản về Ngân Hà
-------------[4] Giải thích hiện tượng thiên văn quan sát được (sao băng, nhật thực, nguyệt thực...)
-------------[5] Thực hành quan sát bầu trời đêm và nhận diện một số chòm sao cơ bản (Bắc Đẩu, Thần Nông...)
-------------[6] Liên hệ ứng dụng công nghệ vũ trụ trong đời sống (GPS, dự báo thời tiết qua vệ tinh, viễn thông...)

----------[4] Ngân Hà
